\sect{Summary \& Research Gap}{summary}

Human-robot interaction has evolved from manual control via teach pendants and vendor-specific programming languages to more flexible interfaces enabled by offline programming, open-source middleware, and API-based integration. NL control methods have progressed from rule-based command parsing to context-aware, multimodal interpretation, and most recently to LLM-driven frameworks capable of translating abstract goals into executable robotic actions. These developments expand the scope of intuitive, high-level human control, reducing the need for low-level programming expertise. In collaborative autonomous environments, combining LLM-based NL control with robotic manipulation, sensing, mobility, and process automation enables fully integrated workflows, allowing users to specify complex objectives in plain language while the system autonomously plans, executes, and adapts tasks across domains ranging from laboratory science to manufacturing and assembly.

Existing NL control frameworks for robotics have demonstrated the ability to map structured or semi-structured commands to robot actions, yet their deployment across collaborative environments remains narrow in scope, often limited to single-purpose systems such as sample preparation, tool transport, or assembly scheduling. Current LLM-integrated solutions primarily handle individual devices or subsystems, lack robust multi-agent or multi-tool orchestration patterns (e.g., manager- or decentralized-style coordination), and struggle with unstructured, conversational inputs that require iterative clarification with the user. Although elements of such agentic self-driving environments may already exist in certain domains, there appears to be no widely established, general-purpose methodology, which enables the systems to orchestrate diverse robotic platforms and instruments within a unified workflow, handle multi-step task sequences, provide real-time feedback. The lack of a clear method for such an integrated, dialogue-driven agent framework suggests a significant research opportunity toward enabling scalable, domain-agnostic autonomous platforms for science, industry, and human-robot collaboration.
